\section{Introducción}

El presente capítulo estima el costo de desarrollo e implementación del sistema,
con el propósito de valorar económicamente el esfuerzo invertido y de determinar
la inversión requerida para su puesta en operación. La estimación se realiza
sobre artefactos verificables ---el software efectivamente construido y el
equipamiento efectivamente empleado--- y no sobre proyecciones abstractas.

El costo se descompone en dos partidas. La primera corresponde al
\textbf{desarrollo de software}, valorado a partir del tamaño funcional del
sistema según el método de puntos de función, contrastado con el tamaño del
código fuente medido directamente. La segunda corresponde al
\textbf{equipamiento} necesario para la operación, considerando que el proyecto
reutiliza la infraestructura de captura existente en el recinto y que, por
consiguiente, la inversión se concentra en la capacidad de cómputo.

Quedan explícitamente fuera del alcance de esta estimación los costos de
operación y mantenimiento posteriores a la puesta en marcha ---consumo
eléctrico, conectividad, soporte y conservación de los equipos---, así como el
costo de la infraestructura de captura preexistente, que no representa un
desembolso atribuible al proyecto.

Todas las valoraciones monetarias de este capítulo se expresan en bolivianos y
se declaran como parámetros de valoración ajustables, dado que dependen de
condiciones de mercado locales que exceden el alcance técnico del proyecto. La
metodología empleada permite recalcular la estimación ante una variación de
dichos parámetros sin alterar el análisis.
